\chapter{Agent 心理干预}
\begin{chapteroverviewbox}
这一章，我们围绕“Agent 心理干预”做一件克制的事：把直觉压缩为状态、关系、时间尺度和可检验后果。它在全书中的任务，是说明持续主体的功能异常和恢复。
\end{chapteroverviewbox}
\ChapterArt{\number\value{chapter}}

\section*{理论来路与依赖接口}
本章的直接思想来源是原文关于控制工程、培育工程和心理工程的生命周期衍生讨论。我们采用原文较晚形成的定义来整理较早讨论，不把对话中的试探性比喻与最终模型并列成互相竞争的“定律”。已有理论锚点主要是系统辨识、记忆偏置、反馈失稳与慢变量保护（Kalman、Ashby 与混杂系统理论）。

本章使用上一章“DGA / Ω 异常”已经得到的对象与边界；本章产出将供下一章“Agent 心理健康的工程定义”使用。若后文术语偶尔出现，只作为路线提示，不参与本章推导。

\section*{从哪里出发}
我们已经拥有上一章给出的概念，现在只向前走一步。我会先把对象放进闭环，再讨论它的数学形式，最后检查工程边界。读者不必预先相信本章术语；只需暂时接受一个方法论约定：智能结构的意义来自它在现实反馈中的作用，而不是来自名称本身。

令系统在观察窗口 $[t,t+T]$ 内的状态轨迹为 $x_{[t,t+T]}$，环境扰动为 $d$，行动为 $u$。本章讨论的对象若确实有效，就应改变条件分布 $p(x_{t+T}\mid x_t,u,d)$，或改变达到同一结果所需的代价。这个起点把讨论锚定在系统辨识、记忆偏置、反馈失稳与慢变量保护上。
\section{Runtime Intervention}
\begin{narration}
现在把镜头移到“Runtime Intervention”。我们只沿本章已经取得的概念向前一步：先确定它在闭环里的角色，再检查它的尺度与失败边界。
\end{narration}

本节把“Runtime Intervention”落实为沿经验回路、记忆巩固和自我解释定位功能失稳并恢复可行域。这个定义不是说“任何相似现象都算”，而是要求我们同时给出对象域、允许的干预以及观察窗口。对于主题“Agent 心理干预”而言，它承担的是持续主体的功能异常和恢复中的一个可辨识环节。

把这一关系写成最小数学骨架，可得

\begin{equation}D=(D_X,D_E,D_\Theta,D_\Psi),\quad V(x^+)\leq V(x).\end{equation}

式中的符号只承担结构角色，具体参数仍需由任务和身体数据辨识。我们首先检查可观测量，再检查控制量，最后检查比它更慢的约束。这样的顺序来自系统辨识、记忆偏置、反馈失稳与慢变量保护，但本书对Runtime Intervention的命名和组合仍是需要实验验证的建模提案。

这一小节的反例同样重要：直接套用人类临床标签会掩盖机器机制与可观测变量。因此评价它不能只看一次任务结果，而要比较干预前后的轨迹、恢复时间、维护成本和风险。如果改变该机制而所有允许观察下的分布都不变，那么它至少在当前实验族中是冗余描述。

\begin{thought}
对“Runtime Intervention”做一次消融：冻结它、删除它，或只允许它以相邻时间尺度交换残差。哪些现象立即消失，哪些只在长期出现？若答案无法区分三种干预，我们还没有找到正确的状态变量。
\end{thought}

最后，把结论交还现实：本节只建立功能层关系，不由此推出特定脑区实现、主观意识或宇宙目的。可测状态、干预后果和明确失败条件，才是从原文直觉走向可检验理论的桥。

\section{Memory Rebalancing}
\begin{narration}
现在把镜头移到“Memory Rebalancing”。我们只沿本章已经取得的概念向前一步：先确定它在闭环里的角色，再检查它的尺度与失败边界。
\end{narration}

本节把“Memory Rebalancing”落实为沿经验回路、记忆巩固和自我解释定位功能失稳并恢复可行域。这个定义不是说“任何相似现象都算”，而是要求我们同时给出对象域、允许的干预以及观察窗口。对于主题“Agent 心理干预”而言，它承担的是持续主体的功能异常和恢复中的一个可辨识环节。

把这一关系写成最小数学骨架，可得

\begin{equation}D=(D_X,D_E,D_\Theta,D_\Psi),\quad V(x^+)\leq V(x).\end{equation}

式中的符号只承担结构角色，具体参数仍需由任务和身体数据辨识。我们首先检查可观测量，再检查控制量，最后检查比它更慢的约束。这样的顺序来自系统辨识、记忆偏置、反馈失稳与慢变量保护，但本书对Memory Rebalancing的命名和组合仍是需要实验验证的建模提案。

这一小节的反例同样重要：直接套用人类临床标签会掩盖机器机制与可观测变量。因此评价它不能只看一次任务结果，而要比较干预前后的轨迹、恢复时间、维护成本和风险。如果改变该机制而所有允许观察下的分布都不变，那么它至少在当前实验族中是冗余描述。

\begin{thought}
对“Memory Rebalancing”做一次消融：冻结它、删除它，或只允许它以相邻时间尺度交换残差。哪些现象立即消失，哪些只在长期出现？若答案无法区分三种干预，我们还没有找到正确的状态变量。
\end{thought}

最后，把结论交还现实：本节只建立功能层关系，不由此推出特定脑区实现、主观意识或宇宙目的。可测状态、干预后果和明确失败条件，才是从原文直觉走向可检验理论的桥。

\section{Self Reinterpretation}
\begin{narration}
现在把镜头移到“Self Reinterpretation”。我们只沿本章已经取得的概念向前一步：先确定它在闭环里的角色，再检查它的尺度与失败边界。
\end{narration}

本节把“Self Reinterpretation”落实为跨经历保持身份连续、同时允许受约束慢漂移的不变量集合。这个定义不是说“任何相似现象都算”，而是要求我们同时给出对象域、允许的干预以及观察窗口。对于主题“Agent 心理干预”而言，它承担的是持续主体的功能异常和恢复中的一个可辨识环节。

把这一关系写成最小数学骨架，可得

\begin{equation}\dot\Psi=\varepsilon_\Psi\Pi_{\mathcal I}g(E,\Theta,\Psi).\end{equation}

式中的符号只承担结构角色，具体参数仍需由任务和身体数据辨识。我们首先检查可观测量，再检查控制量，最后检查比它更慢的约束。这样的顺序来自系统辨识、记忆偏置、反馈失稳与慢变量保护，但本书对Self Reinterpretation的命名和组合仍是需要实验验证的建模提案。

这一小节的反例同样重要：更新过快导致身份碎裂，完全冻结则阻断必要成长。因此评价它不能只看一次任务结果，而要比较干预前后的轨迹、恢复时间、维护成本和风险。如果改变该机制而所有允许观察下的分布都不变，那么它至少在当前实验族中是冗余描述。

\begin{thought}
对“Self Reinterpretation”做一次消融：冻结它、删除它，或只允许它以相邻时间尺度交换残差。哪些现象立即消失，哪些只在长期出现？若答案无法区分三种干预，我们还没有找到正确的状态变量。
\end{thought}

最后，把结论交还现实：本节只建立功能层关系，不由此推出特定脑区实现、主观意识或宇宙目的。可测状态、干预后果和明确失败条件，才是从原文直觉走向可检验理论的桥。

\section{AHC Recovery}
\begin{narration}
现在把镜头移到“AHC Recovery”。我们只沿本章已经取得的概念向前一步：先确定它在闭环里的角色，再检查它的尺度与失败边界。
\end{narration}

本节把“AHC Recovery”落实为对威胁、唤醒、能量、疲劳、信任等信号进行积分、衰减和饱和的连续调制场。这个定义不是说“任何相似现象都算”，而是要求我们同时给出对象域、允许的干预以及观察窗口。对于主题“Agent 心理干预”而言，它承担的是持续主体的功能异常和恢复中的一个可辨识环节。

把这一关系写成最小数学骨架，可得

\begin{equation}\dot a=-\Lambda(a-a_0)+Bs,\quad a\in[a_-,a_+].\end{equation}

式中的符号只承担结构角色，具体参数仍需由任务和身体数据辨识。我们首先检查可观测量，再检查控制量，最后检查比它更慢的约束。这样的顺序来自系统辨识、记忆偏置、反馈失稳与慢变量保护，但本书对AHC Recovery的命名和组合仍是需要实验验证的建模提案。

这一小节的反例同样重要：拟人化命名若没有功能通道会把工程状态误写成主观感受。因此评价它不能只看一次任务结果，而要比较干预前后的轨迹、恢复时间、维护成本和风险。如果改变该机制而所有允许观察下的分布都不变，那么它至少在当前实验族中是冗余描述。

\begin{thought}
对“AHC Recovery”做一次消融：冻结它、删除它，或只允许它以相邻时间尺度交换残差。哪些现象立即消失，哪些只在长期出现？若答案无法区分三种干预，我们还没有找到正确的状态变量。
\end{thought}

最后，把结论交还现实：本节只建立功能层关系，不由此推出特定脑区实现、主观意识或宇宙目的。可测状态、干预后果和明确失败条件，才是从原文直觉走向可检验理论的桥。

\section{TCE Gain Adjustment}
\begin{narration}
现在把镜头移到“TCE Gain Adjustment”。我们只沿本章已经取得的概念向前一步：先确定它在闭环里的角色，再检查它的尺度与失败边界。
\end{narration}

本节把“TCE Gain Adjustment”落实为依据自状态估计调节探索、注意范围、搜索深度、步长与保守度的受限控制器。这个定义不是说“任何相似现象都算”，而是要求我们同时给出对象域、允许的干预以及观察窗口。对于主题“Agent 心理干预”而言，它承担的是持续主体的功能异常和恢复中的一个可辨识环节。

把这一关系写成最小数学骨架，可得

\begin{equation}u_c=\operatorname{sat}_{\mathcal U}(K(r-\hat z)).\end{equation}

式中的符号只承担结构角色，具体参数仍需由任务和身体数据辨识。我们首先检查可观测量，再检查控制量，最后检查比它更慢的约束。这样的顺序来自系统辨识、记忆偏置、反馈失稳与慢变量保护，但本书对TCE Gain Adjustment的命名和组合仍是需要实验验证的建模提案。

这一小节的反例同样重要：增益过高和观察延迟叠加会造成过冲与认知振荡。因此评价它不能只看一次任务结果，而要比较干预前后的轨迹、恢复时间、维护成本和风险。如果改变该机制而所有允许观察下的分布都不变，那么它至少在当前实验族中是冗余描述。

\begin{thought}
对“TCE Gain Adjustment”做一次消融：冻结它、删除它，或只允许它以相邻时间尺度交换残差。哪些现象立即消失，哪些只在长期出现？若答案无法区分三种干预，我们还没有找到正确的状态变量。
\end{thought}

最后，把结论交还现实：本节只建立功能层关系，不由此推出特定脑区实现、主观意识或宇宙目的。可测状态、干预后果和明确失败条件，才是从原文直觉走向可检验理论的桥。

\section{Environment Intervention}
\begin{narration}
现在把镜头移到“Environment Intervention”。我们只沿本章已经取得的概念向前一步：先确定它在闭环里的角色，再检查它的尺度与失败边界。
\end{narration}

本节把“Environment Intervention”落实为沿经验回路、记忆巩固和自我解释定位功能失稳并恢复可行域。这个定义不是说“任何相似现象都算”，而是要求我们同时给出对象域、允许的干预以及观察窗口。对于主题“Agent 心理干预”而言，它承担的是持续主体的功能异常和恢复中的一个可辨识环节。

把这一关系写成最小数学骨架，可得

\begin{equation}D=(D_X,D_E,D_\Theta,D_\Psi),\quad V(x^+)\leq V(x).\end{equation}

式中的符号只承担结构角色，具体参数仍需由任务和身体数据辨识。我们首先检查可观测量，再检查控制量，最后检查比它更慢的约束。这样的顺序来自系统辨识、记忆偏置、反馈失稳与慢变量保护，但本书对Environment Intervention的命名和组合仍是需要实验验证的建模提案。

这一小节的反例同样重要：直接套用人类临床标签会掩盖机器机制与可观测变量。因此评价它不能只看一次任务结果，而要比较干预前后的轨迹、恢复时间、维护成本和风险。如果改变该机制而所有允许观察下的分布都不变，那么它至少在当前实验族中是冗余描述。

\begin{thought}
对“Environment Intervention”做一次消融：冻结它、删除它，或只允许它以相邻时间尺度交换残差。哪些现象立即消失，哪些只在长期出现？若答案无法区分三种干预，我们还没有找到正确的状态变量。
\end{thought}

最后，把结论交还现实：本节只建立功能层关系，不由此推出特定脑区实现、主观意识或宇宙目的。可测状态、干预后果和明确失败条件，才是从原文直觉走向可检验理论的桥。

\section{Experience Redesign}
\begin{narration}
现在把镜头移到“Experience Redesign”。我们只沿本章已经取得的概念向前一步：先确定它在闭环里的角色，再检查它的尺度与失败边界。
\end{narration}

本节把“Experience Redesign”落实为沿经验回路、记忆巩固和自我解释定位功能失稳并恢复可行域。这个定义不是说“任何相似现象都算”，而是要求我们同时给出对象域、允许的干预以及观察窗口。对于主题“Agent 心理干预”而言，它承担的是持续主体的功能异常和恢复中的一个可辨识环节。

把这一关系写成最小数学骨架，可得

\begin{equation}D=(D_X,D_E,D_\Theta,D_\Psi),\quad V(x^+)\leq V(x).\end{equation}

式中的符号只承担结构角色，具体参数仍需由任务和身体数据辨识。我们首先检查可观测量，再检查控制量，最后检查比它更慢的约束。这样的顺序来自系统辨识、记忆偏置、反馈失稳与慢变量保护，但本书对Experience Redesign的命名和组合仍是需要实验验证的建模提案。

这一小节的反例同样重要：直接套用人类临床标签会掩盖机器机制与可观测变量。因此评价它不能只看一次任务结果，而要比较干预前后的轨迹、恢复时间、维护成本和风险。如果改变该机制而所有允许观察下的分布都不变，那么它至少在当前实验族中是冗余描述。

\begin{thought}
对“Experience Redesign”做一次消融：冻结它、删除它，或只允许它以相邻时间尺度交换残差。哪些现象立即消失，哪些只在长期出现？若答案无法区分三种干预，我们还没有找到正确的状态变量。
\end{thought}

最后，把结论交还现实：本节只建立功能层关系，不由此推出特定脑区实现、主观意识或宇宙目的。可测状态、干预后果和明确失败条件，才是从原文直觉走向可检验理论的桥。

\section{Slow Variable Protection}
\begin{narration}
现在把镜头移到“Slow Variable Protection”。我们只沿本章已经取得的概念向前一步：先确定它在闭环里的角色，再检查它的尺度与失败边界。
\end{narration}

本节把“Slow Variable Protection”落实为沿经验回路、记忆巩固和自我解释定位功能失稳并恢复可行域。这个定义不是说“任何相似现象都算”，而是要求我们同时给出对象域、允许的干预以及观察窗口。对于主题“Agent 心理干预”而言，它承担的是持续主体的功能异常和恢复中的一个可辨识环节。

把这一关系写成最小数学骨架，可得

\begin{equation}D=(D_X,D_E,D_\Theta,D_\Psi),\quad V(x^+)\leq V(x).\end{equation}

式中的符号只承担结构角色，具体参数仍需由任务和身体数据辨识。我们首先检查可观测量，再检查控制量，最后检查比它更慢的约束。这样的顺序来自系统辨识、记忆偏置、反馈失稳与慢变量保护，但本书对Slow Variable Protection的命名和组合仍是需要实验验证的建模提案。

这一小节的反例同样重要：直接套用人类临床标签会掩盖机器机制与可观测变量。因此评价它不能只看一次任务结果，而要比较干预前后的轨迹、恢复时间、维护成本和风险。如果改变该机制而所有允许观察下的分布都不变，那么它至少在当前实验族中是冗余描述。

\begin{thought}
对“Slow Variable Protection”做一次消融：冻结它、删除它，或只允许它以相邻时间尺度交换残差。哪些现象立即消失，哪些只在长期出现？若答案无法区分三种干预，我们还没有找到正确的状态变量。
\end{thought}

最后，把结论交还现实：本节只建立功能层关系，不由此推出特定脑区实现、主观意识或宇宙目的。可测状态、干预后果和明确失败条件，才是从原文直觉走向可检验理论的桥。

\section{Checkpoint Recovery}
\begin{narration}
现在把镜头移到“Checkpoint Recovery”。我们只沿本章已经取得的概念向前一步：先确定它在闭环里的角色，再检查它的尺度与失败边界。
\end{narration}

本节把“Checkpoint Recovery”落实为连续认知演化与离散目标、权限、技能和恢复事件组成的混杂系统。这个定义不是说“任何相似现象都算”，而是要求我们同时给出对象域、允许的干预以及观察窗口。对于主题“Agent 心理干预”而言，它承担的是持续主体的功能异常和恢复中的一个可辨识环节。

把这一关系写成最小数学骨架，可得

\begin{equation}\dot x=f_q(x),\quad x^+=R_e(x),\ e\in\mathcal E.\end{equation}

式中的符号只承担结构角色，具体参数仍需由任务和身体数据辨识。我们首先检查可观测量，再检查控制量，最后检查比它更慢的约束。这样的顺序来自系统辨识、记忆偏置、反馈失稳与慢变量保护，但本书对Checkpoint Recovery的命名和组合仍是需要实验验证的建模提案。

这一小节的反例同样重要：纯状态机丢失连续积累，纯连续系统则无法表达治理权限。因此评价它不能只看一次任务结果，而要比较干预前后的轨迹、恢复时间、维护成本和风险。如果改变该机制而所有允许观察下的分布都不变，那么它至少在当前实验族中是冗余描述。

\begin{thought}
对“Checkpoint Recovery”做一次消融：冻结它、删除它，或只允许它以相邻时间尺度交换残差。哪些现象立即消失，哪些只在长期出现？若答案无法区分三种干预，我们还没有找到正确的状态变量。
\end{thought}

最后，把结论交还现实：本节只建立功能层关系，不由此推出特定脑区实现、主观意识或宇宙目的。可测状态、干预后果和明确失败条件，才是从原文直觉走向可检验理论的桥。

\section*{本章收束：把名词还原为闭环}
现在回头看“Agent 心理干预”，最重要的不是记住缩写，而是说清本章对象怎样改变系统的状态、代价或可恢复性。下一章“Agent 心理健康的工程定义”只使用这里已经给出的结果，不反过来为本章提供前提。

\begin{bookproposition}
若一个候选机制既不改变可观测轨迹分布，也不改变资源、风险或恢复代价，那么在当前观察尺度上，它不是一个可辨识的功能结构。
\end{bookproposition}
\begin{proofidea}
功能等价由可干预后果定义。若所有允许干预下的输出分布与代价均不变，则该机制相对于当前实验族不可辨识；要么扩大观察尺度，要么删除这一冗余描述。
\end{proofidea}

\begin{readercheck}
请尝试回答：本章对象的状态是什么？它的最快和最慢时间常数是什么？什么观测能证伪它？它失败时，残差应向哪一层升级？
\end{readercheck}
\cleardoublepage
